% Generated draft: requires mathematical and source review.
\documentclass[11pt]{article}
\usepackage[margin=25mm]{geometry}
\usepackage{fontspec}
\setmainfont{Latin Modern Roman}
\usepackage{amsmath,amssymb,mathtools,unicode-math}
\setmathfont{Latin Modern Math}
\usepackage{enumitem,needspace,xurl,hyperref}
\setlength{\parindent}{0pt}
\setlength{\parskip}{5pt}
\setlength{\emergencystretch}{3em}
\newcommand{\N}{\mathbb{N}}
\newcommand{\Z}{\mathbb{Z}}
\newcommand{\Q}{\mathbb{Q}}
\newcommand{\R}{\mathbb{R}}
\newcommand{\C}{\mathbb{C}}
\newcommand{\F}{\mathbb{F}}
\newcommand{\A}{\mathbb{A}}
\newcommand{\PP}{\mathbb{P}}
\newcommand{\E}{\mathbb{E}}
\newcommand{\eps}{\varepsilon}
\newcommand{\dd}{\,\mathrm{d}}
\begin{document}
\section*{UnsolvedMath ID: 701069}
\textbf{Source problem code:} catalogue entry
\subsection{Definitions and mathematical statement}
A variety \(\mathfrak M\) of groups is a class of groups defined by group identities. A relatively free group \(F_r(\mathfrak M)\) has a free basis \(X_r\) of cardinality \(r\): every map from \(X_r\) to a group in \(\mathfrak M\) extends uniquely to a homomorphism from \(F_r(\mathfrak M)\). An element \(g\in F_r(\mathfrak M)\) is \emph{primitive} if it belongs to some free basis of \(F_r(\mathfrak M)\).

\begin{enumerate}
\item[(a)] Do there exist a variety \(\mathfrak M\), a rank \(r\), an element \(h\in F_r(\mathfrak M)\) that is not primitive, and an injective endomorphism
\[
\alpha:F_r(\mathfrak M)\longrightarrow F_r(\mathfrak M)
\]
such that \(\alpha(h)\) is primitive in \(F_r(\mathfrak M)\)?
\item[(b)] Do there exist a variety \(\mathfrak M\), ranks \(r<n\), and an element \(h\in F_r(\mathfrak M)\) that is not primitive in \(F_r(\mathfrak M)\), but is primitive when regarded as an element of \(F_n(\mathfrak M)\)? Under the usual interpretation of this identification, choose free bases \(X_r\subseteq X_n\) and let
\[
\iota:F_r(\mathfrak M)\hookrightarrow F_n(\mathfrak M)
\]
be the homomorphism induced by the inclusion of bases; the condition is that \(\iota(h)\) is primitive in \(F_n(\mathfrak M)\).
\end{enumerate}

The supplied record does not explicitly specify the permitted range of ranks, including whether they must be finite or positive. It also does not explicitly specify the identification used in part (b); the displayed formulation of that identification uses the conventional inclusion induced by extending a free basis. These qualifications require human review.
\subsection{Short English statement}
E. I. Timoshenko asks whether a nonprimitive element of a relatively free group can become primitive either under an injective endomorphism of that group or upon inclusion into a relatively free group of larger rank in the same variety. The supplied record labels the problem open; this status has not been independently verified. The rank conventions and the intended inclusion in part (b) require clarification.
\subsection{Sources}
\begin{itemize}
\item Supplied local UnsolvedMath registry record: permanent ID \(701069\), exact source code \texttt{catalogue entry}, title \emph{Kourovka Notebook Problem 14.84}. The record attributes the problem to E. I. Timoshenko and gives the proposed year as 1999.
\item \emph{The Kourovka Notebook (Unsolved Problems in Group Theory)}, 14th issue (1999), Problem 14.84; also identified in the supplied background with the 21st issue edition (2026), \texttt{arXiv:1401.0300v48}. These are supplied bibliographic references and have not been independently consulted.
\item The supplied collection description identifies E. I. Khukhro and V. D. Mazurov as editors of the Kourovka Notebook collection covering Issues 1--21 (1965--2026). The local record is marked open and was updated on 2026-10-16; that registry date is not an independently verified literature assessment.
\end{itemize}
\end{document}
